% Requires style.cls, the logo, and references.bib alongside.

\documentclass[twocolumn]{style}
\runninghead{SMILES Summer School Projects Proceedings}
\footertext{\textit{SMILES} 2026}
\addbibresource{references.bib}

\title{Hype Check: Challenging Causal Foundation Models for Uplift}
\logo{SKOLTECH_MACHINE-LEA.png}

\author{
    Vadim Porvatov\textsuperscript{1}
}
\date{
    \textsuperscript{\textbf{1}}
    PJSC Sberbank, Moscow, Russia
}

\begin{document}
\twocolumn[\maketitle]
\small

\begin{abstract}
    \noindent Causal foundation models such as CausalPFN estimate conditional average treatment
    effects (CATE) in context, without per-dataset training, and report strong results on academic
    benchmarks. Uplift modeling, however, is mainly a ranking problem: the goal is to decide which
    units to treat under a budget, not to recover each effect exactly. It is therefore unclear
    whether this advantage holds on real campaigns under a modest compute budget. We first compare
    an untuned CausalPFN with strong tuned baselines on ranking quality, and we test whether effect
    accuracy predicts it. We then identify the regimes where the advantage disappears: data beyond
    the context window, a small control group, and low conversion. Finally, we propose
    CausalPFN-Rank, a training-free adaptation that improves the in-context sample and re-ranks the
    model outputs, and we measure how much ranking quality it recovers. All comparisons are
    relative, use bootstrap confidence intervals, and cover both real randomized trials and
    semi-synthetic data.
\end{abstract}

\noindent\small{\textbf{Index Terms:} Uplift modeling, conditional average treatment effect,
causal foundation models, in-context learning, ranking evaluation.}

\section{Introduction}
Marketing and retention decisions depend on uplift, the incremental effect of an action on an
individual, formalized as the conditional average treatment effect (CATE). In practice the exact
effect value matters less than a good ordering of units by incremental response, so that a limited
budget reaches the most persuadable customers first. Ranking quality is therefore measured with the
Qini coefficient and the area under the uplift curve, not with pointwise error.

Causal foundation models are transformers pre-trained on large collections of simulated causal
models. They estimate CATE in context, with no task-specific training, and report strong benchmark
results~\cite{balazadeh2025causalpfn, robertson2025dopfn, causalfm2025}. The most mature open model,
CausalPFN, ships with pre-trained weights, a simple estimator interface, and calibrated
uncertainty~\cite{balazadeh2025causalpfn}. Recent evidence, though, shows that benchmark accuracy
and real ranking quality can diverge. A large-scale study reports a plain S-Learner topping the Qini
ranking among classical methods~\cite{singh2026criteobench}, and a benchmark on field-collected data
finds most CATE estimates no more accurate than a constant-effect predictor~\cite{yu2024cate}.

We ask whether causal foundation models improve uplift ranking under realistic constraints, where
they break, and whether a lightweight fix repairs them. We propose CausalPFN-Rank, a training-free
adaptation that selects a better in-context sample and re-ranks the model outputs for the ranking
objective. The project has three parts: investigate the model, find where it breaks, and fix it.

\subsection{Hypotheses}
\begin{itemize}
    \item \textit{\textbf{H1:}} Without tuning, CausalPFN matches strong tuned baselines on
    semi-synthetic effect accuracy, but does not beat them on ranking on real campaigns.
    \item \textit{\textbf{H2:}} Ordering methods by effect accuracy (PEHE) agrees only weakly with
    ordering them by uplift (Qini), so the most accurate estimator is often not the one that ranks
    units best.
    \item \textit{\textbf{H3:}} The advantage disappears in three regimes: data beyond the context
    window, a small control group, and low conversion.
    \item \textit{\textbf{H4:}} Engineered context, balancing treated and control and averaging over
    several retrieved samples, recovers much of the ranking lost in the scale and small-control
    regimes, with no training.
    \item \textit{\textbf{H5:}} Re-ranking the outputs improves uplift even when effect accuracy is
    unchanged, narrowing the gap to specialist methods.
\end{itemize}

\subsection{Goals}
\begin{itemize}
    \item \textbf{G1:} Build a reproducible uplift-evaluation harness wrapping CausalPFN and the
    baselines over the dataset pool, and reproduce CausalPFN's own uplift example as an early
    feasibility check.
    \item \textbf{G2:} Map the regimes where the advantage breaks by controlled subsampling,
    reporting uplift metrics with bootstrap confidence intervals, effect accuracy where ground truth
    exists, and the correlation between the two rankings.
    \item \textbf{G3:} Build and ablate CausalPFN-Rank, its context part and its output part, and
    report how much ranking each part recovers, with compute cost.
\end{itemize}

\subsection{Significance of Hypotheses Confirmation}
If validated:
\begin{itemize}
    \item \textit{Practical relevance:} a rule for when an untuned foundation model is enough for
    ranking, and when a tuned classical method or a specialist is worth the effort under a fixed
    budget.
    \item \textit{Methodological value:} evaluation of causal foundation models shifts toward the
    decision objective rather than effect accuracy, and the post-processing recipe applies to any
    in-context causal model.
    \item \textit{Reusable benchmark:} the harness and the map of where the model helps or breaks
    stay useful for a fast-moving area.
\end{itemize}

\subsection{Validation Protocol}
\paragraph{Datasets:} The pool is small and chosen to span the regimes of interest, not to maximize
count. Real randomized data without ground-truth effects: Criteo Uplift and X5 RetailHero at large
scale~\cite{diemert2018criteo}, and Hillstrom, which is small and low-conversion~\cite{hillstrom2008}.
Semi-synthetic data with ground-truth effects: IHDP~\cite{hill2011ihdp} and ACIC~2016~\cite{dorie2019acic}.
We subsample large tables for the scale axis and re-sample for the control-share and conversion axes.

\paragraph{Metrics:}
\begin{itemize}
    \item Ranking: the Qini coefficient and uplift at the top fraction, with bootstrap confidence
    intervals over at least ten seeds.
    \item Accuracy: PEHE against ground-truth effects on semi-synthetic data.
    \item Agreement: Spearman correlation between the PEHE-based and Qini-based rankings of the
    methods.
    \item Efficiency, logged separately: GPU-hours and inference time.
\end{itemize}

\paragraph{Baselines:} The tuned classical core is the S-, T-, X- and DR-Learners on gradient
boosting, plus the Causal Forest~\cite{athey2019grf, battocchi2019econml,
kennedy2023drlearner, kunzel2019metalearners, nie2021rlearner, wager2018causalforest}. We add the
Bayesian Causal Forest and the neural learners DragonNet, TARNet, CFR and
DESCN~\cite{curth2021catenets, hahn2020bcf, shalit2017cfr, shi2019dragonnet,
zhong2022descn}. Two 2026 specialists run on their own regime: the Q-Learner for low conversion
and GP-CATE for a small control group~\cite{fuchs2026qlearner, gpcate2026}. CausalPFN is the
foundation anchor~\cite{balazadeh2025causalpfn}.

\section{Background \& Literature Review}
Under the potential-outcomes framework, the CATE is the expected difference in outcome from treating
an individual with given covariates. In a randomized trial, treatment is independent of covariates,
so the effect is the gap between the treated and control response surfaces and uplift can be
validated directly. The task is to rank units by effect and treat a top fraction under budget, so
ranking quality is summarized by the Qini coefficient and the area under the uplift
curve~\cite{devriendt2018uplift}. This is why a model with worse pointwise accuracy can still rank
better.

Meta-learners reduce effect estimation to supervised learning: the S-, T- and
X-Learners~\cite{kunzel2019metalearners}, the orthogonal R-Learner~\cite{nie2021rlearner}, and doubly
robust DR-Learners~\cite{kennedy2023drlearner}, here all on gradient boosting. Causal Forest and
generalized random forests~\cite{wager2018causalforest, athey2019grf}, and Bayesian Causal
Forest~\cite{hahn2020bcf}, strong on semi-synthetic benchmarks, add mature baselines with
uncertainty. Representation-learning networks form the neural
family~\cite{shalit2017cfr, shi2019dragonnet, zhong2022descn}.

Building on the tabular model TabPFN~\cite{hollmann2025tabpfn}, prior-data-fitted networks have been
specialized for causal effects: CausalPFN~\cite{balazadeh2025causalpfn},
Do-PFN~\cite{robertson2025dopfn}, and the CausalFM framework~\cite{causalfm2025}. CausalPFN is
trained on simulated data-generating processes under ignorability and returns calibrated intervals.
When a table exceeds its context window, it retrieves context rows by a weak-effect rule rather than
reading all rows, and it is not built for instrumental-variable settings. This dependence on a
bounded context, and the gap between the simulated prior and real campaign data, motivate the
breakdown we study.

Two recent non-foundation methods address the hard regimes directly. The Q-Learner writes the
treated-to-control outcome ratio as a product of two odds ratios and is the most consistent method
under low conversion~\cite{fuchs2026qlearner}. GP-CATE models each arm with a Gaussian process and
gives calibrated intervals when the control arm is small, where standard X-Learner intervals
under-cover~\cite{gpcate2026}.

\section{Experimental Setup}
\label{sec:experimental-setup}

This section supplements the original project protocol with the preprocessing and baseline work
completed in the current stage. The broader method pool and robustness experiments specified above
remain part of the planned full study.

\subsection{Datasets}
The baseline study uses four randomized uplift datasets spanning substantially different sample
sizes, treatment balances, and response rates. Criteo is a large-scale advertising experiment and
Hillstrom is an email campaign~\cite{diemert2018criteo, hillstrom2008}. We additionally include the
public X5 RetailHero benchmark and an anonymized LZD campaign dataset. The primary outcome is
conversion for Criteo and Hillstrom and the supplied binary response for LZD and RetailHero. For
Hillstrom, the two email variants are pooled into a single treated arm, so the estimand is the effect
of receiving any campaign email rather than the difference between email creatives. Table
\ref{tab:datasets} summarizes the data used after cleaning.

Before the shared benchmark preprocessing, we use the fixed, audited cleaned versions of all datasets.
For LZD, only the randomized test partition is retained because its original training partition was
created by operational targeting. Near-constant or near-duplicate variables are removed from Criteo
and LZD, while RetailHero variables whose timestamps may extend beyond the campaign are excluded to
avoid post-treatment leakage. These dataset-level decisions are made once, before the experimental
splits, and are therefore identical for every baseline.

\begin{table}[!htpb]
    \caption{Datasets used in the baseline comparison. The feature count is measured after the
    train-fitted transformation.}
    \label{tab:datasets}
    \centering
    \small
    \resizebox{\linewidth}{!}{%
    \begin{tabular}{lrrl}
        \toprule
        \textbf{Dataset} & \textbf{Samples} & \textbf{Features} & \textbf{Outcome} \\
        \midrule
        Criteo & 13,979,592 & 10 & Conversion \\
        Hillstrom & 64,000 & 18 & Conversion \\
        LZD & 181,669 & 64 & Binary response \\
        RetailHero & 200,039 & 5 & Binary response \\
        \bottomrule
    \end{tabular}
    }
\end{table}

\subsection{Preprocessing and Data Splits}
Feature and outcome tables are aligned by unit identifier before any split is constructed. Unit
identifiers, treatment labels, split indicators, treatment timestamps, and known post-treatment
variables are excluded from the covariates. We use a 60/20/20 train--validation--test protocol. For
Hillstrom and RetailHero, splits are stratified by the joint treatment--outcome label. Criteo and
LZD contain repeated covariate vectors; for these datasets we group rows by a hash of the complete
cleaned covariate vector and keep each group in exactly one split. This prevents identical covariates from
appearing in both training and evaluation data while retaining approximately the same split ratios.

All transformations are fitted on the training split only. Missing numerical values are imputed by
the training median and numerical features are standardized. Missing categorical values are imputed
by the training mode and categories are one-hot encoded; categories first observed at evaluation
time are ignored. The fitted transformation is then applied unchanged to validation and test data.
Within a dataset and random seed, every model receives the same sampled rows, split assignments, and
numeric feature matrices.

\subsection{Baseline Estimators and Model Selection}
The T-Learner fits separate treated and control outcome models and ranks units by their predicted
difference. The X-Learner further learns imputed treatment effects in each arm and combines them
using the empirical treatment propensity~\cite{kunzel2019metalearners}. The DR-Learner fits a
gradient-boosted effect model to a cross-fitted doubly robust pseudo-outcome, using the randomized
treatment rate as the propensity~\cite{kennedy2023drlearner}. These three estimators use histogram
gradient boosting for their outcome and effect regressions. DragonNet uses a shared representation,
two potential-outcome heads, and a propensity head~\cite{shi2019dragonnet}. Our implementation uses
the basic joint outcome--propensity objective without targeted regularization; we state this
explicitly to avoid conflating it with the targeted-regularization variant.

For the foundation-model comparison, we use the official CausalPFN implementation and released
pre-trained weights~\cite{balazadeh2025causalpfn}. CausalPFN is not tuned or trained per dataset. It
uses a maximum context length of 4096, a maximum internal query length of 4096, and 1024 retrieved
neighbours per treatment arm; optional temperature calibration is disabled. It receives the same
training split as context and is evaluated on exactly the same test rows as the trained baselines.

Hyperparameters for each dataset--trained-baseline pair are selected by maximizing normalized Qini
AUC on the validation split over 20 configurations. The selected configuration is then
instantiated anew and fitted on the training split; during DragonNet fitting, validation data are
also used to select the early-stopping checkpoint. Each selected model is evaluated once on the
held-out test split. The primary table uses seed 42. To measure sensitivity to splitting and
initialization without repeatedly selecting on validation noise, we freeze these hyperparameters
and repeat the procedure for seeds 0--4. All 80 trained-baseline evaluations in this repeated-seed
matrix are complete. The final study retains the original protocol of at least ten seeds with
bootstrap confidence intervals. No test outcome is used for preprocessing, early stopping, or
hyperparameter selection.

\subsection{Evaluation}
The primary endpoint is normalized Qini AUC, which evaluates whether predicted CATE orders units by
incremental response. We additionally report normalized uplift AUC, the unnormalized AUUC statistic,
uplift among the top 10\% and 20\% of ranked units, and training and prediction time. Because only one
potential outcome is observed for each unit in a real randomized campaign, these metrics evaluate
policy-relevant ranking rather than pointwise CATE error. Pointwise effect error is reserved for the
semi-synthetic experiments with known potential outcomes.

\section{Baseline Comparison}
\label{sec:baseline-results}
Table~\ref{tab:baseline-qini} reports normalized Qini AUC on the held-out seed-42 test splits. There
is no uniformly best estimator. DragonNet is strongest on Criteo and RetailHero, the T-Learner is
best on Hillstrom, and the DR-Learner is best on LZD. The X-Learner is competitive on Criteo but does
not dominate on the smaller campaigns. In particular, DragonNet's strong Criteo result does not
transfer to Hillstrom, where its Qini is negative. These results support treating model family and
campaign regime as part of the comparison rather than assuming that a more complex estimator gives
a consistently better ranking.

\begin{table}[!htpb]
    \caption{Normalized Qini AUC on the held-out test split (seed 42). Higher is better; the best
    result in each dataset is bold.}
    \label{tab:baseline-qini}
    \centering
    \small
    \resizebox{\linewidth}{!}{%
    \begin{tabular}{lrrrr}
        \toprule
        \textbf{Model} & \textbf{Criteo} & \textbf{Hillstrom} & \textbf{LZD} & \textbf{RetailHero} \\
        \midrule
        T-Learner  & 0.476 & \textbf{0.249} & 0.004 & -0.002 \\
        X-Learner  & 0.530 & 0.175 & 0.026 & 0.008 \\
        DR-Learner & 0.496 & 0.238 & \textbf{0.052} & 0.004 \\
        DragonNet  & \textbf{0.545} & -0.042 & 0.023 & \textbf{0.016} \\
        \bottomrule
    \end{tabular}
    }
\end{table}

Table~\ref{tab:baseline-qini-multiseed} summarizes the completed repeated-split evaluation with the
seed-42-selected hyperparameters held fixed. The repeated results change the apparent winners on
Hillstrom and LZD: the DR-Learner has the highest mean on Hillstrom and the X-Learner has the
highest mean on LZD, whereas DragonNet retains the highest mean on Criteo and RetailHero.

\begin{table}[!htpb]
    \caption{Normalized Qini AUC over test splits for seeds 0--4, reported as mean $\pm$ sample
    standard deviation. Hyperparameters are fixed from seed-42 validation selection. Higher is
    better; the highest mean in each dataset is bold.}
    \label{tab:baseline-qini-multiseed}
    \centering
    \small
    \resizebox{\linewidth}{!}{%
    \begin{tabular}{lrrrr}
        \toprule
        \textbf{Model} & \textbf{Criteo} & \textbf{Hillstrom} & \textbf{LZD} & \textbf{RetailHero} \\
        \midrule
        T-Learner  & $0.355 \pm 0.042$ & $-0.003 \pm 0.264$ & $-0.020 \pm 0.031$ & $0.008 \pm 0.003$ \\
        X-Learner  & $0.416 \pm 0.044$ & $0.189 \pm 0.191$ & $\mathbf{0.035 \pm 0.031}$ & $0.017 \pm 0.004$ \\
        DR-Learner & $0.420 \pm 0.057$ & $\mathbf{0.207 \pm 0.109}$ & $0.034 \pm 0.024$ & $0.012 \pm 0.004$ \\
        DragonNet  & $\mathbf{0.435 \pm 0.064}$ & $0.007 \pm 0.218$ & $0.012 \pm 0.035$ & $\mathbf{0.021 \pm 0.003}$ \\
        \bottomrule
    \end{tabular}
    }
\end{table}

The large standard deviations on Hillstrom show that its low event count makes conclusions from a
single split unreliable; in particular, the seed-42 T-Learner advantage does not persist across
seeds. LZD and RetailHero have small absolute Qini values, and the gap between the best two methods
on LZD is negligible relative to split variability. The full-data CausalPFN evaluation is complete
for Hillstrom, LZD, and RetailHero and is running on Criteo, so foundation-model results are kept out
of the table until every dataset--seed cell is available. Bootstrap intervals and the planned
additional seeds remain necessary for final inferential claims.

\section{Method and Timeline}
CausalPFN is used as an in-context estimator: a sample of rows is given as context, and the effect is
read off for new units with no training. We vary three axes: context size on large data, the
control-arm share down to a few-placebo setting, and the base conversion rate. Along each axis we
record where the model's ranking advantage over the classical core disappears.

CausalPFN-Rank keeps the weights frozen and acts in two places. On the input side, it replaces the
default context with a treatment-balanced sample and averages predictions over several retrieved
contexts, compared against that default rather than a random one. On the output side, it adds a light
re-ranking step: a ranking calibrator, a lower-confidence-bound score from the model's uncertainty,
and a ratio score for low conversion. Table~\ref{tab:methods} lists the methods and their roles.

\begin{table}[!htpb]
    \caption{Methods and their role in the study.}
    \label{tab:methods}
    \centering
    \begin{tabularx}{\linewidth}{lX}
        \toprule
        \textbf{Method} & \textbf{Role} \\
        \midrule
        CausalPFN & Foundation anchor \\
        CausalPFN-Rank & Proposed training-free fix \\
        S/T/X/DR + boosting & Tuned classical core \\
        Causal Forest & Forest baseline \\
        BCF & Bayesian baseline \\
        DragonNet, DESCN & Neural baselines \\
        Q-Learner & Specialist: low conversion \\
        GP-CATE & Specialist: small control \\
        \bottomrule
    \end{tabularx}
\end{table}

The project runs over four weeks. Week 1 builds the harness, reproduces CausalPFN's uplift example as
a smoke test, wires up the baselines, and runs the first untuned comparison. Week 2 builds the three
axes and adds the specialists on their regimes. Week 3 implements and ablates the two parts of
CausalPFN-Rank. Week 4 runs the final comparison with confidence intervals and the efficiency
summary, then writes up results and limitations.

\section*{Limitations}
Qini has high variance, so we use bootstrap confidence intervals, at least ten seeds, and a primary
metric fixed in advance. CausalPFN already retrieves context by a weak-effect rule, so the context
part of CausalPFN-Rank is measured against that default, not a random baseline. The model sees only a
bounded context window while gradient boosting trains on all rows; we report this and treat the
context budget as part of the method. Real randomized data allow only ranking evaluation, while
semi-synthetic data give effect accuracy but come from causal models close to the model's training
prior, which is itself the subject of H1. We subsample large tables and train no foundation model.

\printbibliography

\end{document}